\section{AI-Aware Build System}

\label{sec:build}

\subsection{Model Dependency Detection}

The build system automatically scans application code for model dependencies:

\begin{algorithm}[H]
\caption{Model Dependency Detection}
\label{alg:detect}
\begin{algorithmic}[1]
\Require Source code $S$, known model registries $\mathcal{R}$
\State $\text{deps} \gets \text{ParseRequirements}(S)$ \Comment{pip, npm, cargo}
\State $\text{models} \gets \emptyset$
\State \Comment{Pattern 1: HuggingFace model references}
\State $\text{models} \gets \text{models} \cup \text{FindHFModels}(S)$
\State \Comment{Pattern 2: Model download URLs}
\State $\text{models} \gets \text{models} \cup \text{FindModelURLs}(S)$
\State \Comment{Pattern 3: Framework-specific configs}
\State $\text{models} \gets \text{models} \cup \text{FindFrameworkModels}(S)$
\State \Comment{Pattern 4: API client configurations}
\State $\text{apis} \gets \text{FindLLMAPIs}(S)$
\State \Comment{Determine GPU requirements}
\State $\text{gpu} \gets \text{EstimateGPU}(\text{models})$
\State \Return $(\text{models}, \text{apis}, \text{gpu})$
\end{algorithmic}
\end{algorithm}

\subsection{GPU-Optimized Container Images}

For applications requiring local model inference, the build system generates optimized container images:

\begin{enumerate}[leftmargin=1.4em]
    \item \textbf{Base image selection}: Choose from pre-built base images with CUDA, cuDNN, and framework-specific optimizations (PyTorch, TensorFlow, ONNX Runtime, vLLM).
    \item \textbf{Model caching}: Models are stored in a shared cache layer, avoiding redundant downloads across deployments. Cache hit rate: 87\%.
    \item \textbf{Quantization}: Automatically apply GPTQ or AWQ quantization for models that exceed available GPU memory, with quality verification.
    \item \textbf{Multi-stage builds}: Separate build dependencies from runtime, reducing final image size by 40--60\%.
\end{enumerate}

\begin{table}[H]
\centering
\small
\begin{tabular}{lccc}
\toprule
\textbf{App Type} & \textbf{Traditional} & \textbf{Hanzo} & \textbf{Speedup} \\
\midrule
API-only (LLM gateway) & 3.2 min & 1.1 min & 2.9x \\
RAG application & 8.7 min & 2.3 min & 3.8x \\
Local model serving & 42.1 min & 6.8 min & 6.2x \\
Full-stack AI app & 15.4 min & 3.2 min & 4.8x \\
Fine-tuning pipeline & 67.3 min & 12.1 min & 5.6x \\
\bottomrule
\end{tabular}
\caption{Build and deploy times: traditional PaaS vs. Hanzo Platform.}
\label{tab:build_times}
\end{table}

\subsection{Incremental Deployment}

Hanzo Platform supports zero-downtime deployments with AI-specific optimizations:

\begin{algorithm}[H]
\caption{AI-Aware Rolling Deployment}
\label{alg:deploy}
\begin{algorithmic}[1]
\Require New image $I'$, current pods $\mathcal{P}$, model warmup time $T_w$
\State \Comment{Phase 1: Pre-warm new pods}
\State $\mathcal{P}' \gets \text{StartPods}(I', |\mathcal{P}|)$
\For{each $p \in \mathcal{P}'$}
    \State Wait for model loading
    \State Run warmup inference requests
    \State Verify latency $\le$ SLA threshold
\EndFor
\State \Comment{Phase 2: Gradual traffic shift}
\For{$w = 0.1, 0.25, 0.5, 0.75, 1.0$}
    \State Route $w$ fraction of traffic to $\mathcal{P}'$
    \State Monitor error rate and latency for 60s
    \If{error rate $> 1\%$ or P99 latency $> 2\times$ baseline}
        \State Rollback: route all traffic to $\mathcal{P}$
        \State \Return failure
    \EndIf
\EndFor
\State \Comment{Phase 3: Cleanup}
\State Terminate old pods $\mathcal{P}$
\State \Return success
\end{algorithmic}
\end{algorithm}

The model pre-warming phase is critical: without it, cold starts for GPU inference can take 30--120 seconds, causing unacceptable latency spikes during deployment.

